\documentclass{article}
\usepackage[margin=1.5cm]{geometry}
\usepackage{pgfplots}
\usepackage{pgfplotstable}
\pgfplotsset{compat=1.18}
\pagestyle{empty}
\newcommand{\casepath}{CASE_ID}
\begin{document}
\begin{center}
\begin{tikzpicture}
\begin{axis}[width=14cm,height=7cm,xmin=0,xmax=1,ymin=0,
 xlabel={$\theta/\pi$},ylabel={$P(\theta)$ [rad$^{-1}$]},legend pos=north east]
\addplot[blue] table[x=theta_over_pi,y=P] {\casepath/histogram_steps.dat};
\addlegendentry{$P(\theta)$}
\addplot[red,dashed] table[x=theta_over_pi,y=P_reflected] {\casepath/histogram_steps.dat};
\addlegendentry{$P(\pi-\theta)$}
\end{axis}
\end{tikzpicture}

\begin{tikzpicture}
\begin{axis}[width=14cm,height=7cm,xmin=0,xmax=1,ymin=0,ymax=1,
 xlabel={$\theta/\pi$},ylabel={$R(\theta)$},legend pos=north east,
 unbounded coords=jump]
\addplot[purple,mark=*,mark size=1pt] table[x=theta_over_pi,y=R] {\casepath/profile.dat};
\addlegendentry{$R(\theta)$}
\addplot[black,dashed] table[x=theta_over_pi,y=Born] {\casepath/profile.dat};
\addlegendentry{$\cos^2(\theta/2)$}
\end{axis}
\end{tikzpicture}
\end{center}
\clearpage
\begin{center}
\pgfplotstableread{\casepath/nodes.dat}\nodetable
\pgfplotstableread{\casepath/edges.dat}\edgetable
\begin{tikzpicture}[scale=4,detector/.style={circle,draw,fill=white,inner sep=2pt},
 central/.style={circle,draw,fill=orange!30,inner sep=3pt}]
\pgfplotstablegetrowsof{\nodetable}
\pgfmathtruncatemacro{\lastnode}{\pgfplotsretval-1}
\foreach \row in {0,...,\lastnode}{
 \pgfplotstablegetelem{\row}{id}\of\nodetable\pgfmathtruncatemacro{\nid}{\pgfplotsretval}
 \pgfplotstablegetelem{\row}{x}\of\nodetable\edef\nx{\pgfplotsretval}
 \pgfplotstablegetelem{\row}{y}\of\nodetable\edef\ny{\pgfplotsretval}
 \coordinate (v\nid) at (\nx,\ny);
}
\pgfplotstablegetrowsof{\edgetable}
\pgfmathtruncatemacro{\lastedge}{\pgfplotsretval-1}
\foreach \row in {0,...,\lastedge}{
 \pgfplotstablegetelem{\row}{source}\of\edgetable\pgfmathtruncatemacro{\src}{\pgfplotsretval}
 \pgfplotstablegetelem{\row}{target}\of\edgetable\pgfmathtruncatemacro{\dst}{\pgfplotsretval}
 \pgfplotstablegetelem{\row}{kind}\of\edgetable\pgfmathtruncatemacro{\kind}{\pgfplotsretval}
 \ifnum\kind=0 \draw[gray!35,thin] (v\src)--(v\dst);\fi
 \ifnum\kind=1 \draw[blue,thick] (v\src)--(v\dst);\fi
 \ifnum\kind=2 \draw[red,dashed,thick] (v\src)--(v\dst);\fi
}
\foreach \row in {0,...,\lastnode}{
 \pgfplotstablegetelem{\row}{id}\of\nodetable\pgfmathtruncatemacro{\nid}{\pgfplotsretval}
 \pgfplotstablegetelem{\row}{is_central}\of\nodetable\pgfmathtruncatemacro{\centralflag}{\pgfplotsretval}
 \ifnum\centralflag=1 \node[central] at (v\nid) {$q$};
 \else \node[detector,font=\scriptsize] at (v\nid) {\nid};\fi
}
\end{tikzpicture}

Blue: saved detector edges / nearest neighbors. Red dashed: second neighbors.

Gray: central coupling (optional). Labels: zero-based detector IDs.
\end{center}
\end{document}
